\documentclass[12pt]{article}
\usepackage[utf8]{inputenc}
\usepackage[T1]{fontenc}
\usepackage{hyperref}
\title{Xiang Chinese Language Resources for K-12 Curriculum Development}
\author{Compiled by Claude Sonnet 5.5 (claude-sonnet-5-5) for Andr\'as Kornai}
\date{}
\begin{document}
\maketitle

\section{Introduction}
Xiang (hsn) is a Chinese language of Hunan with roughly 38 million speakers, divided into New Xiang and Old Xiang. The search found scholarly and encyclopedic material only; no school textbooks, primers or Xiang-medium curricula, and no dedicated dictionary, were found.

\section{Primary Educational Resources}
None found in search results.

\section{Linguistic and Reference Materials}
\begin{itemize}
\item Yunji Wu, \emph{A Synchronic and Diachronic Study of the Grammar of the Chinese Xiang Dialects} (De Gruyter Mouton). \url{https://www.degruyterbrill.com/document/doi/10.1515/9783110927481/html}
\item Wikipedia: New Xiang (Chang-Yi). \url{https://en.wikipedia.org/wiki/New_Xiang}
\item Wikipedia: Xiangxiang dialect (Central Xiang). \url{https://en.wikipedia.org/wiki/Xiangxiang_dialect}
\item Cambridge article on Changsha Chinese, a Xiang dialect (phonology and tone). \url{https://core-prod.cambridgecore.org/core/product/493730924B25ED441D70D98DBC581513/core-reader}
\item LINGUIST List issue with a review including a chapter on representing Xiang dialects in Chinese characters. \url{https://linguistlist.org/issues/17/107/}
\end{itemize}

\section{Religious and Community Resources}
None found.

\section{Digital and Audio Resources}
None found.

\section{Curriculum Development Context}
Xiang is not known to be taught in schools; instruction is in Standard Mandarin. No standardized Xiang orthography was identified, so character-representation work would be a first step.

\end{document}